% ===== tr-03 块4：第 3 章 pp.32--35（原文页 32--35）=====
% [原文 p.32]
对我们来说另一个重要的事实是：若 $f:Y\to X$ 是态射、$X$ 带有对数结构，并给 $Y$ 配上拉回对数结构，则 $\overline{\mathcal M}_Y=f^{-1}\overline{\mathcal M}_X$。在 $\mathcal M_X$ 由包含 $Y\subseteq X$ 诱导的情形，$\overline{\mathcal M}_X$ 支在 $Y$ 上，于是我们可以把 $\overline{\mathcal M}_X$ 与 $\overline{\mathcal M}_Y$ 等同起来，它们是 $X$ 上除子对数结构及其在 $Y$ 上的限制的幽灵层。
% [Tutor Note: 本段同时出现带与不带两根不同记号：$\overline{\mathcal M}$（带上划线）是幽灵层，$\mathcal M$（无上划线）是真正的对数结构层；原文第一句与末句（推前）用的是后者，切勿一律加上划线]
把这一点与习题 3.8.3 以及对偶相交复形的定义结合起来，我们看到：给定环面退化 $\mathcal X\to D$，对偶相交复形\emph{完全决定}了 $Z$ 之外的幽灵层 $\overline{\mathcal M}_{\mathcal X}=\overline{\mathcal M}_{\mathcal X_0}$。事实上我们采取这样的观点：无论如何，$Z$ 上的对数结构都表现得不特别好，我们总在 $Z$ 上忽略它。事实上，给定 $\mathcal X_0\setminus Z$ 上的对数结构 $\mathcal M_{\mathcal X_0\setminus Z}$，它定义了推前对数结构 $\mathcal M_{\mathcal X_0}:=j_*\mathcal M_{\mathcal X_0\setminus Z}$。由于 $\mathcal X_0$ 是 Cohen--Macaulay 的、从而 $j_*\mathcal O_{\mathcal X_0\setminus Z}=\mathcal O_{\mathcal X_0}$，存在诱导映射 $\alpha:\mathcal M_{\mathcal X_0}\to\mathcal O_{\mathcal X_0}$。因此在下文中，只要我们确定了 $\mathcal X_0\setminus Z$ 上的对数结构，就同样得到 $\mathcal X_0$ 上的对数结构，并且不再关心此对数结构沿 $Z$ 的行为。

这一切给出了倒推、即从 $(B,\mathcal P)$ 走到 $\mathcal X_0^\dagger$ 的必要提示。设给定带奇点的整仿射流形 $B$，配有\emph{环面}多面体分解 $\mathcal P$。在 $\mathcal P$ 的每个顶点 $v$ 处，$\mathcal P$ 局部看起来像一个扇 $\Sigma_v$，它定义环面簇 $X_v$。对每条以 $v$ 与 $w$ 为端点的边 $\omega\in\mathcal P$，$\omega$ 在 $\Sigma_v$ 与 $\Sigma_w$ 两个扇中都定义一条射线，因而给出环面除子 $D_\omega^v\subseteq X_v$、$D_\omega^w\subseteq X_w$。$\mathcal P$ 是环面多面体分解这一条件告诉我们：$D_\omega^v$ 与 $D_\omega^w$ 是同构的环面簇，于是我们可以为每条边 $\omega$ 选取一个环面等变同构 $s_\omega:D_\omega^v\to D_\omega^w$。若选取这些粘合映射使它们在余维 2 分层上满足某个相容性条件（我们把写出这个简单相容性条件留给读者），我们就可以把各个 $X_v$ 粘合起来，得到（一般地）一个代数空间，记作 $X_0(B,\mathcal P,s)$，其中 $s=(s_\omega)$ 是粘合映射的集合。（在 [22] 中，我们以一种略有不同但等价的方式描述粘合数据。）我们称 $s$ 为\emph{闭粘合数据}。我们就是这样构造环面退化的一个潜在中心纤维的。

现在，$X_0(B,\mathcal P,s)$ 不能成为环面退化的中心纤维，除非它带有正确类型的对数结构。这可能不发生的原因很多。首先，若 $s$ 选得不好，$X_0(B,\mathcal P,s)$ 可能存在零维分层，其邻域并不局部平展同构于某个仿射环面簇的环面边界；这是最低限度的先决条件。因此，我们必须把注意力限制在由所谓\emph{开粘合数据}诱导的闭粘合数据上。具体地说，每个极大胞腔 $\sigma\in\mathcal P$ 定义一个仿射环面簇 $U(\sigma)$，它由锥 $C(\sigma)\subseteq M_\R\oplus\R$ 给出（这里把 $\sigma\subseteq M_\R$ 看作格多面体）。设 $V(\sigma)\subseteq U(\sigma)$ 是环面边界。如我们在 [22] 中所证明，$X_0(B,\mathcal P,s)$ 成为环面退化中心纤维的一个必要条件是：它由用平展等价关系除 $\coprod_{\sigma\in\mathcal P_{\max}}V(\sigma)$ 而得。
% [Tutor Note: $\mathcal P_{\max}$ 记 $\mathcal P$ 的极大胞腔之集；$\coprod$ 是不交并，不要与 $\bigcup$ 混；这里的除法"dividing out by an étale equivalence relation"指取该平展等价关系的商（即群胚商），不是代数除法]
换言之，我们是在把各个 $V(\sigma)$ 粘合起来得到一个代数空间，而那些产生形如 $X_0(B,\mathcal P,s)$ 的代数空间的平展等价关系是容易确定的。这在 [22] 第 2 节中有详细展开。由于必须处理代数空间，那里的构造看起来技术上有难度，但基本上是直截了当的。要点是：若 $\sigma_1,\sigma_2\in\mathcal P$ 是两个极大胞腔，$\sigma_1\cap\sigma_2=\tau$，则 $\tau$ 确定了锥 $C(\sigma_1)$ 与 $C(\sigma_2)$ 的面，从而确定开子集 $U_i(\tau)\subseteq U(\sigma_i)$，其环面边界为 $V_i(\tau)\subseteq V(\sigma_i)$。现在一般说来，$U_1(\tau)$ 与
% [原文 p.33]
$U_2(\tau)$ 之间没有\emph{自然}同构：当 $\sigma_1\cap\sigma_2\cap\Gamma\ne\varnothing$ 时这就成了问题，其中 $\Gamma$ 照例是 $B$ 的奇异轨迹。然而关键的是 $V_1(\tau)$ 与 $V_2(\tau)$ 自然同构，于是我们可以选取相容的等变同构，以得到开粘合数据。

选择开粘合数据使我们能定义平展等价关系：我们不过是借助 $V_1(\tau)$ 与 $V_2(\tau)$ 之间选定的同构，把任意两个集合 $V(\sigma_1)$、$V(\sigma_2)$ 粘合起来。开粘合数据的任一选取 $s$ 都以这种方式给出代数空间 $X_0(B,\mathcal P,s)$；而对开粘合数据的任一选取，都有相应的闭粘合数据 $s'$，使得 $X_0(B,\mathcal P,s)\cong X_0(B,\mathcal P,s')$。

使用开粘合数据的好处在于：每个 $\sigma\in\mathcal P_{\max}$ 对应的 $V(\sigma)$ 都带有由除子对数结构 $V(\sigma)\subseteq U(\sigma)$ 诱导的对数结构。遗憾的是，这些对数结构在开粘合映射下并不等同，其原因恰恰是前面提到的 $U_i(\tau)$ 之间缺乏自然同构。不过，这些对数结构的幽灵层是同构的。于是幽灵层 $\overline{\mathcal M}_{V(\sigma)}$ 粘合起来给出幺半群幽灵层 $\overline{\mathcal M}_{X_0(B,\mathcal P,s)}$。把上面所说的总结一下：（这是 [22] 第 2、4 节诸结果的组合）

\begin{theorem}{3.8.4}
给定 $(B,\mathcal P)$，若 $s$ 是闭粘合数据，且 $X_0=X_0(B,\mathcal P,s)$ 是环面退化 $\mathcal X\to D$ 的中心纤维、其对偶相交复形为 $(B,\mathcal P)$，则 $s$ 由开粘合数据诱导，且 $\overline{\mathcal M}_{X_0}|_{X_0\setminus Z}\cong\overline{\mathcal M}_{X_0(B,\mathcal P,s)}|_{X_0\setminus Z}$。
\end{theorem}

这就是我们用组合手段所能达到的极限。下一点是尝试构造 $\mathcal M_{X_0(B,\mathcal P,s)}$。想法是：$\mathcal M_{X_0(B,\mathcal P,s)}$ 是 $\overline{\mathcal M}_{X_0(B,\mathcal P,s)}$ 被 $\mathcal O^\times_{X_0(B,\mathcal P,s)}$ 的一个扩张，所以我们寻找层
\[
\mathcal Ext^1(\overline{\mathcal M}^{\mathrm{gp}}_{X_0(B,\mathcal P,s)},\mathcal O^\times_{X_0(B,\mathcal P,s)})
\]
的某个子层。这里上标 gp 表示幺半群的 Grothendieck 群。$\overline{\mathcal M}^{\mathrm{gp}}_{X_0(B,\mathcal P,s)}$ 被 $\mathcal O^\times_{X_0(B,\mathcal P,s)}$ 的任一扩张都给出一个群层 $\mathcal M^{\mathrm{gp}}_{X_0(B,\mathcal P,s)}$，它满射到 $\overline{\mathcal M}^{\mathrm{gp}}_{X_0(B,\mathcal P,s)}$；而 $\overline{\mathcal M}_{X_0(B,\mathcal P,s)}\subseteq\overline{\mathcal M}^{\mathrm{gp}}_{X_0(B,\mathcal P,s)}$ 的原像是一个幺半群层 $\mathcal M_{X_0(B,\mathcal P,s)}$。
% [Tutor Note: 上式中 gp 是 $\overline{\mathcal M}$ 的上标而非下标，且 $\mathcal Ext$ 用花体（"Ext" 的花体形式）；文字层常把上标 gp 抽成正文，抄写时以页图为准]
当然还需要映射 $\alpha:\mathcal M_{X_0(B,\mathcal P,s)}\to\mathcal O_{X_0(B,\mathcal P,s)}$，这使事情稍复杂一些。长话短说，我们可以辨认出那些给出真正对数结构的扩张所构成的子层。该子层的一个截面决定了 $X_0(B,\mathcal P,s)$ 上一个具有正确幽灵层的对数结构。

不过这并不是我们真正想要的东西。我们真正想要的是 $X_0(B,\mathcal P,s)$ 上的一个对数结构，连同使 $X_0(B,\mathcal P,s)^\dagger\to0^\dagger$ 成为对数光滑的对数态射。（我们稍后会处理坏集 $Z\subseteq\mathcal X_0$ 的问题。）我们把这样的结构称为 $X_0(B,\mathcal P,s)$ 上的一个\emph{对数光滑结构}。结果这类结构由
\[
\mathcal Ext^1(\overline{\mathcal M}^{\mathrm{gp}}_{X_0(B,\mathcal P,s)}/\bar\rho,\mathcal O_X^\times)
\]
的某些截面给出，其中 $\bar\rho$ 是 $\overline{\mathcal M}_{X_0(B,\mathcal P,s)}$ 的典范截面，它在 $\overline{\mathcal M}_{X_0(B,\mathcal P,s),\eta}=\N$ 处的芽为 1，这里 $\eta$ 是 $X_0(B,\mathcal P,s)$ 某个不可约分支的一般点。于是事实上，我们可以辨认出
\[
\mathcal Ext^1(\overline{\mathcal M}^{\mathrm{gp}}_{X_0(B,\mathcal P,s)}/\bar\rho,\mathcal O_X^\times)
\]
的一个子层，记作 $\mathcal{LS}_{X_0(B,\mathcal P,s)}$，其截面决定 $X_0(B,\mathcal P,s)$ 上的一个对数结构\emph{以及}一个对数光滑态射 $X_0(B,\mathcal P,s)^\dagger\to0^\dagger$，即一个对数光滑结构。

[22] 的技术核心是对层 $\mathcal{LS}_{X_0(B,\mathcal P,s)}$ 的显式计算。这在 [22] 的定理 3.22 中局部地完成：该层在 $X_0(B,\mathcal P,s)$ 的（平展）开子集 $V(\sigma)$ 上计算出来；而在 [22] 的定理 3.24 中则是整体地完成。我不打算陈述精确的结果，而是深入一个特殊情形，它能体现该理论最重要的特征。

% [原文 p.34]
\begin{example}{3.8.5}
设 $X_0(B,\mathcal P,s)$ 是正规交叉的，即 $\mathcal P$ 的每个胞腔都仿射同构于标准单形。则我们有局部 $\mathcal T^1$ 层
\[
\mathcal T^1=\mathcal Ext^1_{X_0(B,\mathcal P,s)}(\Omega^1_{X_0(B,\mathcal P,s)/k},\mathcal O_{X_0(B,\mathcal P,s)}).
\]
它是 $S=\operatorname{Sing}(X_0(B,\mathcal P,s))$ 上的线丛。于是可以证明 $\mathcal{LS}_{X_0(B,\mathcal P,s)}$ 就是联系于 $\mathcal T^1$ 的 $\mathcal O_S^\times$-挠子（torsor）。
\end{example}

这使我们回到 Friedman 的 d-半稳定性条件 [10]。带正规交叉的簇是\emph{d-半稳定}的，若 $\mathcal T^1\cong\mathcal O_S$。这样我们就重新得到 Kawamata 与 Namikawa 的结果 [36]：$X_0(B,\mathcal P,s)$ 带有 $0^\dagger$ 上的正规交叉对数结构，当且仅当 $X_0(B,\mathcal P,s)$ 是 d-半稳定的。这当然是因为：联系于 $\mathcal T^1$ 的 $\mathcal O_S^\times$-挠子有截面，当且仅当 $\mathcal T^1\cong\mathcal O_S$。

现在 [22] 的定理 3.24 告诉我们，一般地 $\mathcal{LS}_{X_0(B,\mathcal P,s)}$ 不是平凡的 $\mathcal O_S^\times$-挠子。该层连续地依赖于 $s$，但离散地依赖于 $B$ 的奇点的单值。

若 $\dim B=2$，我们把这个后一要点明确解释一下。$S$ 的不可约分支与 $\mathcal P$ 的一维胞腔一一对应。若 $\omega\in\mathcal P$ 是这样一条边，设它含有 $B$ 的一个奇点，使得 $\Lambda$ 在绕该奇点的某个环路周围的适当基下具有单值
\[
\begin{pmatrix}1&n\\0&1\end{pmatrix}.
\]
则 $\mathcal T^1$ 限制到 $X_0(B,\mathcal P,s)$ 的一维分层 $X_\omega\cong\PP^1$ 上，就是 $\mathcal O_{\PP^1}(n)$。
% [Tutor Note: 这里的 $n$ 就是单值矩阵右上角元素，与 §8 前面"自然数的加法幺半群 $\N$"的记号无关；$n$ 的符号决定 $\mathcal T^1$ 在该 $\PP^1$ 上的度，正文随后正是用 $n<0$ 论证无法光滑化]

为使这个陈述完全准确，需要定义 $n$ 使其与基和环路的选择无关。为此，选取的环路要相对于所选基的定向为逆时针方向，这里的基取自 $\Lambda_b=T_{B,b}$，其中 $b\in B$ 是该环路的基点。

若出现的所有 $n$ 都是正的，则对某些粘合数据 $s$，我们或可期望 $\mathcal T^1$ 有一个只在有限点集 $Z$ 上消失的截面 $t$。若 $Z$ 不含环面分层（即不含三重点），则我们得到 $X_0(B,\mathcal P,s)\setminus Z$ 上所需类型的对数结构，从而通过推前得到 $X_0(B,\mathcal P,s)$ 上的对数结构（在 $Z$ 之外对数光滑）。于是我们有

\begin{proposition}{3.8.6}
在本例的情形下，若 $\dim B=2$ 且
\[
t\in\Gamma(X_0(B,\mathcal P,s),\mathcal T^1)
\]
是在不含三重点的有限集 $Z$ 上消失的截面，则存在 $X_0(B,\mathcal P,s)$ 的光滑化 $\mathcal X\to D$，使得 $\mathcal X$ 的奇异轨迹是 $Z\subseteq\mathcal X_0=X_0(B,\mathcal P,s)$，且诱导的对数态射 $\mathcal X_0^\dagger\to0^\dagger$ 与由 $t$ 决定的 $X_0(B,\mathcal P,s)^\dagger\to0^\dagger$ 相合。
\end{proposition}

这个证明是 Friedman 或 Kawamata--Namikawa 结果的一个相当简单的应用。然而为了应用这些结果，我们需要处理奇异集 $Z$。做法是：把 $X_0(B,\mathcal P,s)$ 正规化，在每个 $Z$ 中点的原像中各爆破一个点，然后沿导子轨迹的真变换重新粘合。这样得到一个 Friedman 意义下的 d-半稳定簇，或一个对数光滑概形，它随后可以被光滑化。（这样的做法在高维似乎很困难。）

另一方面，若 $B$ 的某个奇点处 $n<0$，我们就会遇到问题，而且此时确实不存在 $X_0(B,\mathcal P,s)$ 的光滑化。这不应令人惊讶，理由如下。若 $n=-1$，结果是我们将不得不在这样的奇点上方添加一种奇怪的 $I_1$ 纤维来紧化环面纤维化 $X(B_0)$。$I_1$ 纤维是一个浸入球面，而对通常的 $I_1$ 纤维，
% [原文 p.35]
两片在该纤维奇点处的相交重数为 $+1$。然而当单值由 $n=-1$ 给出时，相交重数为 $-1$。这对特殊拉格朗日 $T^2$-纤维化是不会出现的，因此在这种情况下构造不出光滑化并不奇怪。\proofend

若 $\dim B=2$ 且对 $B$ 上所有奇点都有 $n>0$，则我们称 $B$ 是\emph{正的}。可以把"正"这一概念推广到带多面体分解的高维 $B$，见 [22] 的定义 1.54。$B$ 的正性是 $X_0(B,\mathcal P,s)$ 能作为环面退化中心纤维出现的必要条件。第 4 节的所有例子都是正的；这事实上可由自反多面体的凸性推出，而正性条件可以看作一种凸性陈述。

现在回到一般情形，不对 $B$ 的维数或 $\mathcal P$ 的胞腔形状作任何限制。由 [22] 的定理 3.24 可知，$\mathcal{LS}_{X_0(B,\mathcal P,s)}$ 是某个凝聚层------我们称之为 $\mathcal{LS}^+_{pre,X_0(B,\mathcal P,s)}$------的集的子层。该层是直和
\[
\bigoplus_{\substack{\omega\in\mathcal P\\ \dim\omega=1}}\mathcal N_\omega,
\]
其中 $\mathcal N_\omega$ 是 $X_0(B,\mathcal P,s)$ 中对应于 $\omega$ 的环面分层上的线丛。此外，与二维正规交叉情形一样，若 $B$ 是正的，则 $\mathcal N_\omega$ 是半丰沛线丛。
% [Tutor Note: $\mathcal{LS}^+_{pre,X_0(B,\mathcal P,s)}$ 中 $+$ 是上标、$\mathrm{pre}$ 是下标的首段，不要写成 $(\mathcal{LS}_{pre})^+$ 之类；直和的指标是两条限制 $\omega\in\mathcal P$ 与 $\dim\omega=1$ 同时成立]

截面 $t\in\Gamma(X_0(B,\mathcal P,s),\mathcal{LS}^+_{pre,X_0(B,\mathcal P,s)})$ 若在 $t$ 的零点集 $Z$ 之外是 $\mathcal{LS}_{X_0(B,\mathcal P,s)}$ 的截面，则它决定了 $X_0(B,\mathcal P,s)\setminus Z$ 上一个对数光滑结构。特别地，若 $Z$ 不含任何环面分层，则一切顺利。于是我们得到在 $Z$ 之外对数光滑的对数态射 $X_0(B,\mathcal P,s)^\dagger\to0^\dagger$。我们把这样的结构称为一个\emph{对数 Calabi--Yau 空间}。

我们来回顾一下：给定数据
\begin{itemize}
\item $s$：开粘合数据；
\item $t\in\Gamma(X_0(B,\mathcal P,s),\mathcal{LS}^+_{pre,X_0(B,\mathcal P,s)})$，其中 $t$ 在 $X_0(B,\mathcal P,s)\setminus Z$ 上是 $\mathcal{LS}_{X_0(B,\mathcal P,s)}$ 的截面，而 $Z$ 是不含 $X_0(B,\mathcal P,s)$ 的任何环面分层的某集合；
\end{itemize}
我们便得到 $X_0(B,\mathcal P,s)^\dagger\to0^\dagger$。

反之，我们在 [22] 中证明：若 $\mathcal X\to D$ 是环面退化，则 $\mathcal X_0^\dagger\to0^\dagger$ 可由对偶相交复形 $(B,\mathcal P)$ 经某种数据 $s$、$t$ 的选择按上述方式得到。为完成这幅图景，还有待回答

\begin{question}{3.8.7}
（重构问题，版本 III）设 $(B,\mathcal P)$ 是正的。
\begin{enumerate}[label=(\arabic*)]
\item 哪些 $s$ 与 $t$ 的选取能产生对数 Calabi--Yau 空间？
\item 给定 $X_0(B,\mathcal P,s)^\dagger\to0^\dagger$，它何时能作为某个环面退化 $\mathcal X\to D$ 的中心纤维出现？
\end{enumerate}
\end{question}

按我们目前的勾勒，这个问题是基本重构问题 3.7.11 的精细版本。

数据 $s$ 与 $t$ 的选取决定了由给定对偶相交复形产生的对数 Calabi--Yau 空间的模空间。到目前为止我们甚至还没有断言这个模空间非空；而对一般的 $(B,\mathcal P)$，我并不知道它是否非空，尽管当 $\dim B=2$ 或 $3$ 时它确实非空。但无论如何，人们希望对这个模空间有更明确的描述。一般地，该模空间是一个概形，但它可能是奇异的（[22] 的例 4.28 给出了一个例子）。要解决这个问题，还需要一些附加假设。为引出必要的假设，让我们回到第 1 节，那里我们引入了单性（simplicity）的概念。我们看到：镜像对称的基本拓扑，只在纤维化简单时才成立。所以或许我们应当期望：当前的构造在具有单性时表现得更好。
% [说明：以上一段起于原文 p.35、止于 p.36 前三句；为保持段落完整，整段留在本文件，p.36 的页标记放在 tr03e 开头（该页下一段开始处）]
